\documentclass{leveldoc}
\DocUses{foundations,slingshot,fast_lane,first_magnet,calutron}
\DocFingerprint{7dc161c85022571e}

\begin{document}

\section{What the level asks}\label{sec:asks}

The game's text:
\leveldescription

Six positive ions leave the source at $(\val{xZero}, \val{yZero})$ with $T_0 = \val{Tzero}$:
three masses, $1$, $2$ and $4$, each at $\pm\val{angDeg}$ degrees from $+x$. Every flight must
first pass through the gate, the box $\val{gx0} \le x \le \val{gx1}$,
$\val{gy0} \le y \le \val{gy1}$ (\physref{6.2}: a detector entered with the gate still missing
ends the flight as a skipped gate), and then enter the detector of its mass at
$x = \val{xDet}$ (the table in \S\ref{sec:hand}); there are no acceptance windows. The player
places up to five charges and four magnets from the listed values in the region
$\val{regX0} \le x \le \val{regX1}$, $\val{regY0} \le y \le \val{regY1}$. A flight ends when the
ion touches a charge or a magnet, leaves the world, enters its detector, or reaches
$t_{\max} = \val{tmax}$ (\xsecref{foundations}{sec:board}).

This is the magnetic module's finale: an instrument in stages. The reference's two negative
charges on the axis gather each mass's diverging pair through the gate, as a lens does
(\levelref{fast_lane}'s impulse of a passing charge, \xeqref{fast_lane}{eq:impulse}); its three
magnets then sort the masses by their rigidity $p/q$
(\xeqref{calutron}{eq:curvature}), the lightest bent highest. This document integrates the six
flights numerically.

\section{Setup}\label{sec:setup}

The level file, field by field (the appendix reproduces it):
\setuptable

In symbols: every ion has charge $q = \val{q}$ and radius $0$; $c = \val{c}$. The reference's
elements (charges of radius $\val{R}$, magnets of radius $\val{Rm}$, field units, \physref{2.2}):
\begin{center}\small\val{elements}\end{center}
The world is $-1 \le x \le \val{xWorld}$, $-1 \le y \le \val{yWorld}$.

\begin{figure}[htbp]
\centering
\begin{tikzpicture}
\begin{axis}[axis equal image, width=0.95\linewidth, xmin=-1.6, xmax=41.6, ymin=-1.6, ymax=21.6,
  xtick={0,5,...,40}, ytick={0,5,...,20}, tick label style={font=\small},
  xlabel={$x$}, ylabel={$y$}, label style={font=\small}]
  \draw[black!8] (0,0) grid (40,20);
  \draw[black!45, dashed] (-1,-1) rectangle (\rawval{xWorld},\rawval{yWorld});
  \draw[violet!70!black, thick] (\rawval{gx0},\rawval{gy0}) rectangle (\rawval{gx1},\rawval{gy1});
  \draw[green!40!black] (37,15) rectangle (40,20);
  \draw[green!40!black] (37,8) rectangle (40,13);
  \draw[green!40!black] (37,4) rectangle (40,7);
  \addplot[blue!70!black] table {patha.dat};
  \addplot[blue!70!black, dashed] table {pathb.dat};
  \addplot[red!70!black] table {pathc.dat};
  \addplot[red!70!black, dashed] table {pathd.dat};
  \addplot[orange!80!black] table {pathe.dat};
  \addplot[orange!80!black, dashed] table {pathf.dat};
  \fill (\rawval{xZero},\rawval{yZero}) circle[radius=0.15];
\end{axis}
\end{tikzpicture}
\caption{The gate (violet), the three detectors (green; the level file's) and the six flights:
mass 1 (blue), 2 (red) and 4 (orange), each at $-\val{angDeg}^\circ$ (solid) and
$+\val{angDeg}^\circ$ (dashed). The elements are listed in the table above.}
\label{fig:board}
\end{figure}

\section{Principles and model}\label{sec:model}

\begin{principles}
  \item The Lorentz force in the relativistic equation of motion \xeqref{foundations}{eq:eom};
    the charges' field \xeqref{foundations}{eq:coulomb}; the magnets' in-plane field
    \xeqref{first_magnet}{eq:dipole-plane}; the energy \xeqref{foundations}{eq:energy}.
  \item The board's rules and the gate (\xsecref{foundations}{sec:board}, \physref{6.2}); the
    classical Runge--Kutta method, cubic Hermite interpolation and the error estimate by halving
    the step (\xsecref{foundations}{sec:numerics}), as in \xsecref{slingshot}{sec:method}.
\end{principles}

The assumptions are First magnet's (\xsecref{first_magnet}{sec:model}): the elements are held in
place, the ions are points without moments, and their radiation is not taken from them.

\section{Derivation}\label{sec:derivation}

\subsection{The stages}\label{sec:stages}

A negative charge on the axis attracts each ray towards the axis by an impulse that grows as the
ray passes closer (\xeqref{fast_lane}{eq:impulse}: the angle $2\kappa/(b\,pv)$ for a straight pass
at the distance $b$); placed where the two rays of a mass are on either side, it turns them
towards each other, and the gate on the axis selects the rays it has gathered. The ions then have
the same energy but momenta in the ratio $1 : \sqrt2 : 2$, and in the magnets' field each path
curves at the rate $qB_z/p$ (\xeqref{calutron}{eq:curvature}): the light ions are bent most. The
charges' field also acts beyond the gate, so the exit angles are not the magnetic integral alone;
the energy \xeqref{foundations}{eq:energy} is the check.

\subsection{The method}\label{sec:method}

The six flights are integrated by \xeqref{foundations}{eq:rk4}; the gate is passed when a sample
of a step's cubic Hermite interpolant lies in it, the next event located as in
\xsecref{slingshot}{sec:method}. The step $h = \val{H}$ and its half are both run, and the finer
flight is reported: the step-halving estimate $\abs{y_h - y_{h/2}}/15$ is the error of the finer
flight (of the coarser it is sixteen times that). The largest estimate is $\val{errMax}$, ion~2's,
whose path passes close to an element; the energy drifts by at most $\val{wDevMax}$.

\section{The reference solution by hand}\label{sec:hand}

\begin{procedure}
  \item\label{st:march} \emph{The march.} Advance each ion by \xeqref{foundations}{eq:rk4} in the
    elements' fields at the step $h/2$, check the gate, and locate its entry:
    \begin{center}\small\val{results}\end{center}
    Every ion passes the gate and enters its detector through its left face, $x = \val{xDet}$,
    inside the detector's span.
\end{procedure}

\begin{result}[The reference flights]
All six ions pass the gate and enter the detectors of their masses at the points of the table.
The level is solved.
\end{result}

\section{Comparison with the game}\label{sec:comparison}

\comparisontable

Each flight's tolerance is the largest of $10^{-9}$, ten times its own step estimate and the
game's own preview--verification gap, fixed before comparing. The largest difference, ion~2's
ordinate, is the size of its step estimate: at its close pass both integrations work harder (the
game's preview and verification differ there by $\val{gapTwo}$); smaller steps converge onto the
game's verified value.

\section{Other solutions}\label{sec:others}

On the empty board, the second placement in \code{placements.json}, the ions fly straight and
diverge; most leave through the right edge, and ion~5 enters a detector without having passed the
gate (\val{bareOutFive}). A solution needs a stage that gathers each pair through the gate and one
that sorts the masses, which no single element does. A map of solutions would take the game's
verdicts, which this document does not reproduce.

\section{What the game computes}\label{sec:engine}

The game integrates the equations of motion with its adaptive eighth-order Runge--Kutta method
(\physref{5}), passes gates by its certified crossing search and tracks their margins
(\physref{6.2}), and verifies the outcome at a tighter tolerance (\physref{7}); the differences
in the comparison are the two integrations' errors.

\section{Further reading}\label{sec:reading}

The stages' physics: \levelref{fast_lane} (the impulse of a passing charge),
\levelref{first_magnet} (a magnet's field), \levelref{calutron} (rigidity); motion in a magnetic
field: \Jackson{§12.2}.

\levelfileappendix

\end{document}
